تقترح هذه الورقة سير عمل لإخفاء المعلومات اللغوي الواعي بالسياق وفق نموذج \textbf{الصندوق الأسود} في النقاشات المتسلسلة، إذ تحمل الحمولة السرية عبر \textit{قناة اختيار} بدلًا من التلاعب على مستوى الرموز. يعيد المرسل بناء مجموعة وثائق سياقية مشتركة، ويختار هدف الرد والاستطرادة الدلالية، ثم يستخدم نموذج لغة كبيرًا عشوائيًا لصياغة هذا الاختيار على شكل رد مرئي عادي. يعيد المستقبل بناء السياق نفسه، ويسترجع الاختيارات المحددة، ويفك ترميز الحمولة من التعليق المنشور. صُمِّمت الطريقة كي تبقى طبيعية داخل المحادثة، وتتجنب علامات التحكم المرئية، وتدعم الاسترجاع اعتمادًا على الافتراضات المشتركة والنافذة الزمنية المتفق عليها فقط. نقيّم النهج باستخدام مقاييس الطبيعية والتباعد التوزيعي (\EN{GPT-2 PPL}، \EN{KL/JSD} على أحاديات الكلمات) على 183 عينة ناجحة من أصل 200 عينة في تشغيل نظيف، ونُورد إخفاقات التوليد وفك الترميز بصورة منفصلة.
